\section{An elementary polynomial-time primality test}
\label{app:primality}

The factorization procedure needs an exact test for its prime leaves and
for the prime factors appearing in its certificates.  We give the required
test here, including its correctness proof. It is a convenient variant of
the polynomial-congruence method of Agrawal, Kayal, and
Saxena~\cite[Section~4]{AKS};
the auxiliary prime and constants used below are proved directly.

\begin{lemma}\label{lem:primality}
There is a uniform deterministic algorithm which, given integers
$n\ge128$ and $2\le m<2^n$, decides whether $m$ is prime using a number of
bit operations polynomial in $n$.
\end{lemma}

\begin{proof}
We first describe the algorithm.
\begin{enumerate}
\item Reject $m$ if it is a perfect power $b^e$ with integers $b,e\ge2$.
\item Find a prime $s\le n^8$ such that
\begin{equation}\label{eq:primality-order}
 \gcd(s,m)=1,
 \qquad \ord_s(m)>n^2.
\end{equation}
\item If $m\le8s$, decide primality by trial division and stop.
Otherwise reject if any integer from $2$ through $8s$ divides $m$.
\item Accept if and only if all the congruences
\begin{equation}\label{eq:primality-congruences}
 (X+a)^m\equiv X^m+a\pmod{m,\,X^s-1},
 \qquad 1\le a\le8s,
\end{equation}
hold.
\end{enumerate}
We prove that the search in the second step succeeds, that the algorithm
is correct, and that its running time is polynomial.

\paragraph{Existence of $s$.}
Let $R=2\lfloor n^8/2\rfloor$ and let $P_R$ be the product of the primes
at most $R$.  The central binomial coefficient $C=\binom{R}{R/2}$ satisfies
$C\ge2^R/(R+1)$.  For every prime $\ell$, the factorial valuation formula gives
\[
 v_\ell(C)=\sum_{j\ge1}
 \left(\left\lfloor\frac{R}{\ell^j}\right\rfloor
       -2\left\lfloor\frac{R/2}{\ell^j}\right\rfloor\right)
 \le\lfloor\log_\ell R\rfloor.
\]
Each summand is zero or one, and it vanishes when $\ell^j>R$.
Consequently, if $\pi(R)$ denotes the number of primes at most $R$, then
\[
 \log_2 C\le\pi(R)\log_2 R
 \le(\log_2 P_R)\log_2 R.
\]
Thus
\begin{equation}\label{eq:primality-prime-product}
 \log_2 P_R
 \ge\frac{R-\log_2(R+1)}{\log_2 R}>n^6.
\end{equation}
For the last inequality, $8\log_2 n+1<n$ for $n\ge128$;
the difference is positive at $128$ and increasing thereafter.
Since $R\ge n^8-1$ and $R+1\le2n^8$, the displayed quotient is larger than
$(n^8-1-n)/n$, which is larger than $n^6$.
On the other hand,
\[
 \log_2\left(m\prod_{j=1}^{n^2}(m^j-1)\right)
 <n\left(1+\frac{n^2(n^2+1)}2\right)<n^6.
\]
It follows from \eqref{eq:primality-prime-product} that some prime
$s\le R$ divides neither $m$ nor any $m^j-1$ with $1\le j\le n^2$.
This is precisely \eqref{eq:primality-order}.

\paragraph{Correctness.}
Primes pass the perfect-power and trial-division checks.  If $m$ is prime,
the binomial theorem in characteristic $m$ gives
$(X+a)^m=X^m+a^m=X^m+a$, so they also pass
\eqref{eq:primality-congruences}.

Suppose, conversely, that a composite $m$ is accepted.  It is not a perfect
power, and every prime $p$ dividing it satisfies $p>8s$.  Choose one such
$p$.  In an algebraic extension of $\mathbb F_p$, choose a root $w$ of
$1+X+\cdots+X^{s-1}$.  Its value at $1$ is $s\ne0$ in $\mathbb F_p$,
so $w\ne1$.  Since $s$ is prime, $w$ has multiplicative order exactly $s$.

For any product $f$ of the tested linear factors $X+a$, reduction of
\eqref{eq:primality-congruences} modulo $p$ gives
$f(X)^m\equiv f(X^m)\pmod{X^s-1}$.  Frobenius gives the same identity with
exponent $p$.  These identities compose: substitution $X\mapsto X^b$
preserves the ideal $(X^s-1)$, so identities for exponents $b,c$ imply the
identity for $bc$.  Hence, for all nonnegative integers $u,v$,
\begin{equation}\label{eq:primality-propagation}
 f(w)^k=f(w^k),\qquad k=m^u p^v.
\end{equation}

Let $H$ be the subgroup of $(\Z/s\Z)^\times$ generated by the residues
of $m$ and $p$, and put $h=|H|$.  By \eqref{eq:primality-order},
\[
 n^2<h\le s-1.
\]
The nonnegative powers $m^u p^v$ already range over $H$, because inverses
in a finite group are nonnegative powers.  We now obtain incompatible
lower and upper bounds on the number of values $f(w)$.
These bounds adapt the root-counting arguments
in~\cite[Lemmas~4.7--4.8 and their footnotes]{AKS}, which credit the
lower-bound argument to Lenstra and independently Macaj, and the
upper-bound formulation to Kalai, Sahai, and Sudan.

For the lower bound, take all products of $h-1$ distinct factors from
$X+1,\ldots,X+8s$.  The factors are pairwise distinct over $\mathbb F_p$
because $p>8s$, so different subsets give different monic polynomials.
If two such polynomials $f_1,f_2$ had the same value at $w$, then
\eqref{eq:primality-propagation} would give
$f_1(w^k)=f_2(w^k)$ for every $k\bmod s\in H$.  This would give $h$
distinct roots of the nonzero polynomial $f_1-f_2$, whose degree is at
most $h-1$.  Thus all these values are distinct, and their number is at least
\begin{equation}\label{eq:primality-lower}
 \binom{8s}{h-1}
 \ge\left(\frac{8s}{h-1}\right)^{h-1}
 \ge8^{h-1}.
\end{equation}
Here $\binom{M}{t}\ge(M/t)^t$ follows by comparing each factor
$(M-i)/(t-i)$ with $M/t$.  No division by $f(w)$ has been used.  Indeed
$f(w)=0$ is itself impossible, since \eqref{eq:primality-propagation}
would give $h$ roots of the nonzero polynomial $f$, of degree $h-1$.

For the upper bound, put $q=\lfloor\sqrt h\rfloor$.  The $(q+1)^2>h$
pairs $0\le u,v\le q$ give two distinct pairs for which the corresponding
integers $k_1=m^u p^v$ and $k_2=m^{u'}p^{v'}$ agree modulo $s$.
They are different as integers: since $m$ is composite and not a perfect
power, it has a prime divisor $p'\ne p$.  Equality of $k_1$ and $k_2$
would force $u=u'$ by taking $p'$-adic valuations, and then $v=v'$.
Now $w^{k_1}=w^{k_2}$, so \eqref{eq:primality-propagation} makes every
value counted in \eqref{eq:primality-lower} a root of the nonzero polynomial
$T^{k_1}-T^{k_2}$.  Its degree is at most
\[
 \max(k_1,k_2)\le m^{2q}
 <2^{2nq}\le2^{2n\sqrt h}<2^{2h}.
\]
But \eqref{eq:primality-lower} is at least
$8^{h-1}=2^{3h-3}>2^{2h}$, since $h>n^2\ge128^2$.
This contradiction proves correctness.

\paragraph{Running time.}
Every exponent in a representation $m=b^e$ with $b\ge2$ satisfies $e<n$,
so binary search for integer roots with exponents $2,\ldots,n$ implements
the first step in polynomial time.  The search for $s$ uses trial division
to recognize primes up to $n^8$, followed by gcd and successive modular
powers to test \eqref{eq:primality-order}.  These operations are polynomial
in $n$ and do not use the primality test under construction.  The trial
divisions involving $m$ are also bounded by a polynomial in $n$.
Finally, elements of $(\Z/m\Z)[X]/(X^s-1)$ are represented by $s$
coefficients of at most $n$ bits.  Schoolbook cyclic multiplication and
repeated squaring test each of the $8s$ congruences using polynomially many
bit operations, since $s\le n^8$ and the exponent $m$ has at most $n$ bits.
The prime factor $p$, the extension field, and the exponents in the root
count are used only in the correctness proof and need not be computed.
All loops have uniform polynomial bounds, completing the proof.
\end{proof}
